\section{Actual bounds for the retained four-window costs}

Continue with (HT1)--(HT18), now with $k\geq3$.  All unadorned
source norms and quotient volumes in this section are those of the original
$m_{h,k}(u)\,du$, with $g=2\xi$ and the complete hypothetical quartet.
Write $\mathcal L=\mathcal L_{h,k}$ and retain $q=\ell_k(k+1)^2$.
The threshold-four lower bound is inherited from
\emph{Consecutive identities and the first window}, (CJ.9)--(CJ.19);
it is already explicitly present in the recovered transcript at lines
4202--4228.  It is not a new claim of the present continuation.

\subsection{Every term of the inherited lower bound}

Put $d_j=V_j/V_{j-1}$ for $j\geq q$.  For the two original windows
$I_-=(q-1,2q-1)$ and $I_+=(q,2q)$, retain the contraction penalties
\[
 \begin{split}
 P_-&=-\log(1-d_{2q-1})-2\sum_{j=q}^{2q-2}\log(1-d_j),\\
 P_+&=-\log(1-d_q)-\log(1-d_{2q})
       -2\sum_{j=q+1}^{2q-1}\log(1-d_j).
 \end{split}\tag{HC1}
\]
The source identities prove $0<d_j<1$ in these ranges.  The real phase
$\phi_N$ is defined by the retained source pairing
$b_N^*G_Nb_{N+1}=i\omega_N\phi_N$.  The first window uses its original
boundary identity; no inverse at degree $q-2$ occurs.  Define weights
\[
 w_{q-1}=w_{2q-1}=1,\qquad w_N=2\quad(q\leq N\leq2q-2),
 \qquad\sum_{N=q-1}^{2q-1}w_N=2q.
\]
Adding the two exact source products gives
\[
 \mathcal B_k:=\log\frac{V_{q-1}V_q}{V_{2q-1}V_{2q}}
 =\sum_{N=q-1}^{2q-1}w_N\log(\varepsilon_N^2+\phi_N^2)
 -\log\frac{\omega_{2q-1}\omega_{2q}}{\omega_{q-1}\omega_q}
 +P_-+P_+.
 \tag{HC2}
\]
Here every $\varepsilon_N\geq\mathcal L$, including the first degree.

For clarity, retain all constants in the inherited comparable-window
bound.  With $\alpha=\pi/2$, $B=42+8m$, and
$\vartheta_h=\int_{-1}^{1}w_h$, put
\[
 \begin{gathered}
 t_h=\min(1,\vartheta_h),\quad c_0=\min(1,2c_h/3),\quad
 \ell_h=\tfrac12 e^{-\alpha}\sqrt{c_0t_h}\,2^{-B/2},\\
 X=\max(1,M_h(b),M_h(-b)),\quad U_h=4\sqrt X/b,\\
 C_h^{\rm win}=27\max(1,U_h)^3\max(1,\ell_h^{-1})^2,
 \qquad D_h=\frac{\delta}{2C_h^{\rm win}}.
 \end{gathered}\tag{HC3}
\]
Here $0<b<\pi/2$, and $c_h>0$ is the original proved convolution
constant.  The source proves
$(\omega_{n+r}/\omega_n)^{1/(2r)}\leq C_h^{\rm win}n$ for
$n\geq k\geq3$, $n/2\leq r\leq2n$.  Thus it applies to both literal
windows, and $0<D_h<1/108$.

The exact residual in (HC2) has the following five nonnegative pieces:
\[
 \begin{split}
 \mathcal R_k={}&P_-+P_+
 +\sum_Nw_N\log(1+\phi_N^2/\mathcal L^2)\\
 &+\sum_Nw_N\log\frac{\varepsilon_N^2+\phi_N^2}
                              {\mathcal L^2+\phi_N^2}\\
 &+4q\log(C_h^{\rm win}q)
       -\log\frac{\omega_{2q-1}\omega_{2q}}{\omega_{q-1}\omega_q}
 +4q\log\frac{2\mathcal L}{\delta kq}.
 \end{split}\tag{HC4}
\]
Both sums use exactly $q-1\leq N\leq2q-1$.  Nonnegativity follows,
respectively, from $0<d_j<1$, the squared phases, (HT9), the two
comparable-window estimates, and $\mathcal L/q\geq\delta k/2$.
Expanding each logarithm in (HC4), with $D_h=\delta/(2C_h^{\rm win})$,
proves the exact identity
\[
 \boxed{\mathcal B_k=4q\log(D_hk)+\mathcal R_k,
 \qquad\mathcal R_k\geq0.}\tag{HC5}
\]
This retains, rather than discards, all the contraction, phase, norm,
and trace savings in the original lower bound.

\subsection{A proved upper bound for these same four original volumes}

The original convolution estimate and Laplace moments are
\[
 m_{h,k}(u)\geq c_h\vartheta_h^{k-3}
 e^{-\alpha(|u|+k-3)}(1+|u|+k-3)^{-B},\qquad
 M_h(\pm b)=\int_{\mathbb R}e^{\pm bu}w_h(u)\,du.
\]
Write $\mu_h=\int_{\mathbb R}w_h(u)\,du$ for the original source mass.
Put
\[
 \begin{gathered}
 R_0=\sqrt{\delta^2+\gamma^2},\quad
 D=\max(1,2/b,2R_0),\quad T=Dq,\quad r=kR_0/T\leq1/2,\\
 a_k(T)=c_h\vartheta_h^{k-3}e^{-\alpha(T+k-3)}
 (1+T+k-3)^{-B},\\
 M_k=M_h(b)^k+M_h(-b)^k,\qquad
 U_k(T)=\mu_h^k+\frac{(2q-2)!}{(bT)^{2q-2}}M_k,\\
 A_N=\sqrt{\frac{2N+1}{2}}\frac{4^{N+1}-1}{3},\qquad
 B_q(r)=\left(\frac{1+r}{1-r}\right)^q,\qquad
 C_N=\frac{U_k(T)A_N^2B_q(r)^2}{T a_k(T)}.
 \end{gathered}\tag{HC6}
\]
The capital $D$ in the interval scale is distinct from the inherited
lower constant $D_h$ and from the source differential.  All masses,
coordinates, and root multiplicities remain as in (HT1)--(HT7).
One has the actual finite Hermitian comparisons
\[
 G_{q-1}\preceq C_{2q-1}G_{2q-1},\qquad
 G_q\preceq G_{q-1}\preceq C_{2q}G_{2q}.
\]
Consequently, with the exact ratio
\[
 \rho_q=\frac{4q+1}{4q-1}
 \left(\frac{4^{2q+1}-1}{4^{2q}-1}\right)^2,
\]
\[
 \boxed{0\leq\mathcal B_k
 \leq q\log C_{2q-1}+q\log C_{2q}
 =2q\log C_{2q-1}+q\log\rho_q.}\tag{HC7}
\]

\paragraph{Proof of the degree extension and its maps.}
Use the exact substitution $f(x)=P(c+iTx)$ with inverse
$P(S)=f((S-c)/(iT))$.  Its source measure is $T m_{h,k}(Tx)\,dx$,
with mass $\mu_h^k$, and its monic relation is
$\chi_T(x)=(iT)^{-q}\chi(c+iTx)$.  The roots $z_1,\ldots,z_q$,
listed with their repetitions, have modulus at most $r$.
The raw jet of order $d$ acquires exactly $(iT)^d$; the quotient
determinant acquires $T^{-q(q-1)}$.  This common factor cancels in
each displayed quotient-volume ratio.

The Legendre expansion on $[-1,1]$ gives
$\|f\|_{{\rm coeff},1}\leq A_N\|f\|_{L^2([-1,1])}$ for
$\deg f\leq N$.  Indeed Rodrigues' polynomial has coefficient norm
at most $\binom{2j}{j}\leq4^j$ and squared integral $2/(2j+1)$.
Pairing the expansion with this polynomial bounds its coefficient by
$\sqrt{(2j+1)/2}\|f\|_2$; summing the geometric series gives exactly
$A_N$ in (HC6).

Write $\prod_a(1-z_at)^{-1}=\sum_{j\geq0}h_jt^j$.
Geometric-series multiplication gives
$|h_j|\leq\binom{q+j-1}{j}r^j$ and
$\|\chi_T\|_{{\rm coeff},1}\leq(1+r)^q$.
For every integer $s\geq0$, monic division of $x^{q+s}$ has quotient
$\sum_{j=0}^sh_jx^{s-j}$.  The inverse-series identity cancels all
other terms of degree at least $q$.  Removing the leading coefficient
one therefore gives
\[
 \|\operatorname{rem}_{\chi_T}x^{q+s}\|_{{\rm coeff},1}
 \leq(1+r)^q\sum_{j=0}^s\binom{q+j-1}{j}r^j-1
 \leq B_q(r).
\]
The bound is valid also for $s=q$, the extra degree required here.
At degrees below $q$ the remainder is the original monomial, whose
coefficient norm is one.  Linearity gives the same remainder-map
bound for every degree at most $2q$.

For a remainder $v$ of degree below $q$,
\[
 \int |v(u/T)|^2m_{h,k}(u)\,du
 \leq U_k(T)\|v\|_{{\rm coeff},1}^{2}.
\]
This follows by bounding each monomial by
$\max(1,|u/T|^{q-1})$ and using
$\int u^{2q-2}m_{h,k}\leq(2q-2)!b^{-(2q-2)}M_k$.
Restriction of the original norm to $[-T,T]$ gives the lower bound
$T a_k(T)\|f\|_{L^2([-1,1])}^2$.  Combining these three estimates proves
\[
 \|\operatorname{rem}_{\chi}P\|_{H_{q-1}}^2
 \leq C_N\|P\|_{H_N}^2\qquad(N=2q-1,2q).
\]
The map is the literal remainder representative of the original jet.
Apply it to the canonical least lift of $x\in E$ at degree $N$.
Its norm at degree $q-1$ is at least $x^*G_{q-1}x$; this proves the
Hermitian comparisons preceding (HC7).  Inclusion of polynomial degrees
gives $G_q\preceq G_{q-1}$ and monotonicity of all volumes.  Taking
determinants proves (HC7).  Applying the same remainder inequality to
any nonzero polynomial below degree $q$ proves $C_N\geq1$ without an
additional assumption.\hfill$\square$

For an entirely scalar upper bound retain the original constants
\[
 K=\max(1,\vartheta_h^{-1}),\qquad
 C_{\rm env}=\max(1,2/(3c_hD)),\quad a=\alpha D+\log256,
\]
and the original expression (AU.31)
\[
 \mathcal U_k=a q^2+
 \left(\log(XK)+\alpha+\frac{16R_0}{3D}\right)kq
 +Bq\log(1+(D+1)q)+q\log C_{\rm env}.
\]
Then the proved four-window upper bound is
\[
 \boxed{\mathcal B_k\leq\widehat{\mathcal U}_k
 :=2\mathcal U_k+q\log20.}\tag{HC8}
\]
To verify the constants, $bT\geq2q$ implies the factorial ratio in
(HC6) is at most one; also $\mu_h\leq X$ and $M_k\leq2X^k$, so
$U_k(T)\leq3X^k$.  Directly from (HC6),
\[
 A_{2q-1}^2\leq(2q/9)256^q,\qquad
 A_{2q}^2\leq(40q/9)256^q.
\]
The logarithmic derivative $2/(1-r^2)\leq8/3$ on $[0,1/2]$ gives
$\log B_q(r)^2\leq16kR_0/(3D)$.  Substitution into $C_N$ yields
$q\log C_{2q-1}\leq\mathcal U_k$ and
$q\log C_{2q}\leq\mathcal U_k+q\log20$.
Adding proves (HC8).  This argument does not use the false inequality
$\rho_2<20$; in the actual quartet domain $q\geq16$ anyway.

\subsection{The exact compatible interval for the signed arithmetic correction}

Let $V_N^\Gamma$ denote the inherited Gamma-reference quotient volume
for the \emph{same} original polynomial $\chi$, and define
\[
 T_N=V_N/V_N^\Gamma,\qquad
 \mathcal B_k^\Gamma=\log\frac{V_{q-1}^\Gamma V_q^\Gamma}
 {V_{2q-1}^\Gamma V_{2q}^\Gamma},\qquad
 \Delta_k^\Gamma=\log\frac{T_{q-1}T_q}{T_{2q-1}T_{2q}}.
\]
Cancellation of the four literal volume factors gives
$\mathcal B_k=\mathcal B_k^\Gamma+\Delta_k^\Gamma$.
No auxiliary polynomial is introduced and no sign is imposed on
$\Delta_k^\Gamma$.  From (HC5) and (HC8) the actual retained residuals
must belong to the following explicit intervals. Let $\mathfrak C_*$
be (AT18a) on the same arithmetic and tensor Gamma moment matrices.
Its proof uses the original full packet and applies to exactly these
four determinants, so $|\Delta_k^\Gamma|\le J_*\le\mathfrak C_*$.
Write $T_k=4q\log(D_hk)$ for this displayed scalar only:
\[
 \boxed{
 \max\{0,-T_k,\mathfrak b_*^--T_k\}
 \leq\mathcal R_k
 \leq\min\{\widehat{\mathcal U}_k,\mathfrak b_*^+\}-T_k,}
 \tag{HC9}
\]
\[
 \boxed{
 \max\{\mathfrak d_*^-,\max\{0,4q\log(D_hk)\}-\mathcal B_k^\Gamma\}
 \leq\Delta_k^\Gamma
 \leq\min\{\mathfrak d_*^+,\widehat{\mathcal U}_k-\mathcal B_k^\Gamma\}.}
 \tag{HC10}
\]
The new entries follow by substituting
$\mathcal B_k=\mathcal B_k^\Gamma+\Delta_k^\Gamma$ and
$\mathcal R_k=\mathcal B_k-T_k$ into (AT18)--(AT19).
Here $\mathfrak b_*^\pm,\mathfrak d_*^\pm$ are the
simultaneous original Gamma endpoints (AT21e); its complete actual
moment proof supplies both the additive and nonlinear coordinate.
Their intersection width has not been evaluated for an actual
offcritical packet. Removing the simultaneous (AT21e) endpoint
entries $\mathfrak b_*^\pm$ and $\mathfrak d_*^\pm$ gives
the preceding coarse intervals, whose width is strictly positive
for every $k\geq3$.
Indeed all nonquadratic terms of $\mathcal U_k$ are nonnegative,
$a>1$, $q\geq(k+1)^2>2k$, and $D_h<1$.  Thus
\[
 \widehat{\mathcal U}_k\geq2a q^2
 >4qk\geq\max\{0,4q\log(D_hk)\}.
\]
Moreover, keeping the fixed packet constants, each lower-order term
divided by $q^2$ tends to zero, so
\[
 \lim_{k\to\infty}
 \frac{\widehat{\mathcal U}_k-\max\{0,4q\log(D_hk)\}}{q^2}
 =2(\alpha D+\log256)>0.\tag{HC11}
\]
This proves numerical compatibility of those coarse bound functions.
It does not prove compatibility of the refined (HC9)--(HC10),
which retain the source-dependent (AT21e) endpoints computed
from $J_*$ and $I_*$.  It does not construct a packet realizing an arbitrarily
chosen point of either interval, and it does not assert the existence
of the hypothetical zero.  In particular the interval may lie entirely
at negative values of $\Delta_k^\Gamma$; determinant-ratio positivity
does not imply positivity of its signed four-window logarithm.


\paragraph{The constructed signed refinement on this identical source.}
The exact coordinate map to (ACM1)--(ACM5) is $y=u$, $x=t$,
with $S=k/2+iu$ and every ascending coefficient unchanged.
The reference density (ACM2) is exactly the tensor Gamma density
of (AT4), including its mass $(\sqrt{2\pi})^k$; the arithmetic
density is the original $m_{h,k}=w_h^{*k}$. The relation map in
the common source is $I_NB_N^{\rm AW}$ in both presentations.
Consequently the source Grams, their quotient minima, and all
four signed determinants agree. In particular $F(0)=\mathcal B_k^\Gamma$
and $F(1)=\mathcal B_k$ for the $F$ of (ACM14)--(ACM16).

Here is the exact dictionary for all pointwise controls:
\[
 (S_{\rm AW},T_q,S_{\rm SP},S_q,A_q)
 =(S_{\rm spec},S_{\rm tr},S_{\rm ov},S_{\rm HS},S_{\rm ang}),
 \qquad \mathscr H_{2q-1}=H_{2q-1}.
\]
The original crossed-pair proof (AT21b) gives
$S_{\rm tr}\le S_{\rm sin}\le S_{\rm ang}$, and the original
overlap proof gives $S_{\rm tr}\le S_{\rm ov}$. Therefore the
full pointwise minima in (ACM11a) and (AT21a)--(AT21c) are
identical, as are those minima divided by
$(2q-1)\sqrt{1-e^{-F(x)/(2q-1)}}$. Thus $J_q=J_*$ and $I_q=I_*$
on this actual path. These equalities retain both complete angle
lists, including their designated and additional zero entries.
They compare the calculated expressions without removing any
source vector or changing either endpoint form.

For every finite-series degree $p\ge0$, let
$j_{p,\Gamma},e_{p,\Gamma}$ be the exact signed center and residual
constructed in (ACM15) on this Gamma-to-arithmetic path.
Their actual finite geometric identity and centered
Hilbert--Schmidt proof give
$j_{p,\Gamma}-e_{p,\Gamma}\le\Delta_k^\Gamma
\le j_{p,\Gamma}+e_{p,\Gamma}$.
The explicit residual in (ACM15) tends to zero for this fixed
source pair, with every eigenvalue and the factor $8q-6$ retained.
Substituting the preceding dictionary into (ACM16) gives exactly
\[
\begin{aligned}
 \beta_{p,\Gamma}^-&=
 \max\{\mathfrak b_*^-,
                  \mathcal B_k^\Gamma+j_{p,\Gamma}-e_{p,\Gamma}\},\\
 \beta_{p,\Gamma}^+&=
 \min\{\mathfrak b_*^+,
                  \mathcal B_k^\Gamma+j_{p,\Gamma}+e_{p,\Gamma}\},\\
 \mathfrak b_*^-&\le\beta_{p,\Gamma}^-
 \le\mathcal B_k\le\beta_{p,\Gamma}^+\le\mathfrak b_*^+ .
\end{aligned}\tag{HC9a}
\]
Every lower entry is at most the same actual $\mathcal B_k$ and
every upper entry is at least it, proving the displayed
intersection and its order. Their signed entries also prove
$0\le\beta_{p,\Gamma}^+-\beta_{p,\Gamma}^-\le2e_{p,\Gamma}$.
Indeed the upper endpoint is at most
$\mathcal B_k^\Gamma+j_{p,\Gamma}+e_{p,\Gamma}$ and the lower
endpoint is at least
$\mathcal B_k^\Gamma+j_{p,\Gamma}-e_{p,\Gamma}$.
Thus both endpoints converge to the actual $\mathcal B_k$ on
this fixed source as $p\to\infty$; no nesting in $p$ is required.
Inserting (HC5) and subtracting the
same original Gamma value yields the stronger intervals
\[
\begin{aligned}
 \max\{0,-T_k,\beta_{p,\Gamma}^--T_k\}
 &\le\mathcal R_k
 \le\min\{\widehat{\mathcal U}_k,\beta_{p,\Gamma}^+\}-T_k,\\
 \max\{0,T_k,\beta_{p,\Gamma}^-\}-\mathcal B_k^\Gamma
 &\le\Delta_k^\Gamma\\
 &\le\min\{\widehat{\mathcal U}_k,\beta_{p,\Gamma}^+\}
                         -\mathcal B_k^\Gamma .
\end{aligned}\tag{HC10a}
\]
The first line uses $\mathcal R_k=\mathcal B_k-T_k\ge0$,
and the last two use $\Delta_k^\Gamma=\mathcal B_k-\mathcal B_k^\Gamma$.
The nesting in (HC9a), together with
$\mathfrak d_*^\pm=\mathfrak b_*^\pm-\mathcal B_k^\Gamma$,
proves that these intervals lie inside (HC9)--(HC10).
Those original bounds and their full proofs are consequently
retained corollaries. Removing all source-dependent endpoint
entries gives the original coarse intervals in (HC11); its
positive width is not an evaluation of the refined width.
All signed centers here are the specified exact integrals.
Any numerical evaluation must retain its integration error;
fixed-source convergence in $p$ supplies no uniform estimate in $k$.
\paragraph{The graph calculation carried back into these arithmetic intervals.}
On the additional full-quartet domain $\gamma>2$ of (ACM17)--(ACM27),
retain $k\geq3$ as in this section. The coefficients, polynomial
$\chi$, original $v_h=2\xi/h$ and quotient maps are identified
there literally with (HT1)--(HT3). The graph computation therefore
controls these same arithmetic quantities by
$\mathcal B_k=\widehat{\mathcal B}_k-\Gamma_k^{\rm gr}$.
For any specified integers $p,s\geq0$, define
\[
 \begin{aligned}
 b_{p,s}^{\rm gr,-}
 &=\max\{\beta_{p,\Gamma}^-,
                     \widehat{\mathcal B}_k-g_{s,\rm gr}^+\},\\
 b_{p,s}^{\rm gr,+}
 &=\min\{\beta_{p,\Gamma}^+,
                     \widehat{\mathcal B}_k-g_{s,\rm gr}^-\}.
 \end{aligned}
 \tag{HC9g}
\]
The $\beta_{p,\Gamma}^{\pm}$ retain the original
Gamma-to-arithmetic path; $g_{s,\rm gr}^{\pm}$ are the separately
computed arithmetic-to-graph endpoints (ACM27). Subtracting the
two bounds on $\Gamma_k^{\rm gr}$ from this same
$\widehat{\mathcal B}_k$ gives
\[
 \beta_{p,\Gamma}^-\leq b_{p,s}^{\rm gr,-}
 \leq\mathcal B_k\leq b_{p,s}^{\rm gr,+}
 \leq\beta_{p,\Gamma}^+.
\]
Inserting this exact intersection at the use sites (HC9)--(HC10)
and keeping their original $T_k=4q\log(D_hk)$ proves
\[
 \begin{aligned}
 \max\{0,-T_k,b_{p,s}^{\rm gr,-}-T_k\}
 &\leq\mathcal R_k
 \leq\min\{\widehat{\mathcal U}_k,b_{p,s}^{\rm gr,+}\}-T_k,\\
 \max\{0,T_k,b_{p,s}^{\rm gr,-}\}-\mathcal B_k^\Gamma
 &\leq\Delta_k^\Gamma\\
 &\leq\min\{\widehat{\mathcal U}_k,b_{p,s}^{\rm gr,+}\}
                                      -\mathcal B_k^\Gamma.
 \end{aligned}
 \tag{HC10g}
\]
Indeed $\mathcal R_k=\mathcal B_k-T_k\geq0$ by (HC5),
$\mathcal B_k\geq0$ by (HC7), and
$\Delta_k^\Gamma=\mathcal B_k-\mathcal B_k^\Gamma$ by its
definition. Subtracting these unchanged constants gives every
entry and proves nesting into the preceding intervals. The
additional interval width obeys
\[
 0\leq b_{p,s}^{\rm gr,+}-b_{p,s}^{\rm gr,-}
       \leq\min\{2e_{p,\Gamma},2e_{s,\rm gr}\}.
 \tag{HC10h}
\]
For the first bound, (ACM16) places the two Gamma endpoints
inside the interval centered at
$\mathcal B_k^\Gamma+j_{p,\Gamma}$ with radius $e_{p,\Gamma}$.
For the second, (ACM27) places the graph-corrected endpoints
inside the interval centered at
$\widehat{\mathcal B}_k-j_{s,\rm gr}$ with radius $e_{s,\rm gr}$.
Taking an intersection cannot increase either width. These are
actual finite-source enclosures. Their fixed-source residual
limits do not estimate the separate limit in the growing tensor
order. All bounds preceding this paragraph retain their original
domain $\gamma>0$; (HC9g)--(HC10h) assert the additional calculated
refinement on the explicitly stated graph-source domain.


\paragraph{The original relation-row enclosure at these same use sites.}
Retain every equation (HC9g)--(HC10h), and keep its domain
$k\geq3$, $\gamma>2$, $0<\delta<1/2$, $m\geq1$.
Equations (ACM28)--(ACM31) compute the additional lower and upper
values $L_k^{\rm rel},U_{k,2q}^{\rm rel}$ from the full monic
relation rows and the original finite arithmetic graph Grams.
Here $d_j^{\rm rel}$ is a relation polynomial as defined there;
the original $d_j=V_j/V_{j-1}$ of (HC1) remains its stated scalar.
Likewise the form $T_N^{\rm rel}:=f_N^*f_N$ below retains its
superscript and does not rename either $T_N=V_N/V_N^\Gamma$ or
the scalar $T_k=4q\log(D_hk)$. With
$g_s^{\pm}=g_{s,\rm gr}^{\pm}$ exactly as in (ACM27), put
\[
\begin{aligned}
 b_{p,s}^{\rm rel,-}
   &=\max\{\beta_{p,\Gamma}^-,L_k^{\rm rel}-g_{s,\rm gr}^+\},\\
 b_{p,s}^{\rm rel,+}
   &=\min\{\beta_{p,\Gamma}^+,U_{k,2q}^{\rm rel}-g_{s,\rm gr}^-\},\\
 \beta_{p,\Gamma}^-&\leq b_{p,s}^{\rm rel,-}
 \leq b_{p,s}^{\rm gr,-}\leq\mathcal B_k
 \leq b_{p,s}^{\rm gr,+}\leq b_{p,s}^{\rm rel,+}
 \leq\beta_{p,\Gamma}^+.
\end{aligned}\tag{HC9r}
\]
The coefficient identity (ACM17)--(ACM20) identifies the graph
and arithmetic remainder frames by the literal identity, including
the full primary nilpotents. Consequently (ACM31) bounds the same
$\mathcal B_k$ as (HC9a). Taking their interval intersection proves
the displayed relation enclosure. For its comparison with (HC9g),
$L_k^{\rm rel}\leq\widehat{\mathcal B}_k$ gives
$\max\{\beta^-,L_k^{\rm rel}-g^+\}
\leq\max\{\beta^-,\widehat{\mathcal B}_k-g^+\}$, and
$\widehat{\mathcal B}_k\leq U_{k,2q}^{\rm rel}$ gives the
corresponding order for the two upper minima. These two elementary
monotonicities prove every remaining order in (HC9r). The exact
graph-volume interval is contained in the relation-row interval;
its intersection with the latter is itself. Thus (HC9g) retains its
full exact input, while (HC9r) gives the separately calculated
finite row enclosure for that input.

At the two receiving residuals define
\[
\begin{aligned}
 L_{k;p,s}^{\rm ar,rel}&=\max\{0,T_k,b_{p,s}^{\rm rel,-}\},\\
 U_{k;p,s}^{\rm ar,rel}&=\min\{\widehat{\mathcal U}_k,
                                              b_{p,s}^{\rm rel,+}\},\\
 L_{k;p,s}^{\rm ar,rel}-T_k
 &\leq\mathcal R_k\leq U_{k;p,s}^{\rm ar,rel}-T_k,\\
 L_{k;p,s}^{\rm ar,rel}-\mathcal B_k^\Gamma
 &\leq\Delta_k^\Gamma
                  \leq U_{k;p,s}^{\rm ar,rel}-\mathcal B_k^\Gamma.
\end{aligned}\tag{HC10r}
\]
To prove this complete substitution, (HC5) gives
$\mathcal B_k\geq T_k$ and
$\mathcal R_k=\mathcal B_k-T_k\geq0$, (HC7)--(HC8) give
$0\leq\mathcal B_k\leq\widehat{\mathcal U}_k$, and (HC9r)
gives its third lower and second upper entries. Their intersection
contains this same scalar. The two maps are the translations
$x\mapsto x-T_k$ and $x\mapsto x-\mathcal B_k^\Gamma$,
with inverse maps adding the displayed unchanged constant.
They are order-preserving bijections of the real line, so they
preserve interval widths and give the two complete receiving
intervals. The first lower endpoint is exactly
$\max\{0,-T_k,b_{p,s}^{\rm rel,-}-T_k\}$.
No sign is assigned to $g_{s,\rm gr}^-$, to $T_k$, or to
$\Delta_k^\Gamma$ before these explicit operations.
Monotonicity of the maximum and minimum also proves that (HC10g)
is contained in (HC10r), which is contained in (HC10a).

The relation-row width has the precise bound
\[
\begin{aligned}
0\leq b_{p,s}^{\rm rel,+}-b_{p,s}^{\rm rel,-}
&\leq\min\{2e_{p,\Gamma},\,
 U_{k,2q}^{\rm rel}-L_k^{\rm rel}
                      +g_{s,\rm gr}^+-g_{s,\rm gr}^-\}\\
&\leq\min\{2e_{p,\Gamma},\,
 U_{k,2q}^{\rm rel}-L_k^{\rm rel}+2e_{s,\rm gr}\}.
\end{aligned}\tag{HC10s}
\]
The first lower sign follows because the interval contains
$\mathcal B_k$. Its Gamma width is at most $2e_{p,\Gamma}$
by (HC9a). Its other interval has lower endpoint $L_k^{\rm rel}-g^+$
and upper endpoint $U_{k,2q}^{\rm rel}-g^-$, so direct subtraction
gives the second entry. Equation (ACM27) bounds $g^+-g^-$ by
$2e_{s,\rm gr}$. Taking an intersection cannot increase either
of these two widths, which proves the stated order. The complete
negative mixed squared norm in (ACM30) reduces the row upper
comparison relative to its $K_{\rm low}^{-1}$ term. It does
not alter the exact-graph nesting just proved. All fixed-source
limits in $p$ and $s$ retain their stated fixed $h,k$ domain.


\paragraph{The full correlation calculation in the retained signed intervals.}
Keep the exact row subsets $J(j)$ of (ACM35), on the same
receiving domain as (HC9r). Their calculated bounds now give
\[
\begin{aligned}
 b_{p,s}^{\rm corr,-}
 &=\max\{\beta_{p,\Gamma}^-,L_k^{\rm corr}(J)-g_{s,\rm gr}^+\},\\
 b_{p,s}^{\rm corr,+}
 &=\min\{\beta_{p,\Gamma}^+,U_k^{\rm corr}(J)-g_{s,\rm gr}^-\},\\
 b_{p,s}^{\rm rel,-}&\leq b_{p,s}^{\rm corr,-}
 \leq b_{p,s}^{\rm gr,-}\leq\mathcal B_k
 \leq b_{p,s}^{\rm gr,+}\leq b_{p,s}^{\rm corr,+}
 \leq b_{p,s}^{\rm rel,+}.
\end{aligned}\tag{HC9rc}
\]
The exact identity $\mathcal B_k=\widehat{\mathcal B}_k-\Gamma_k^{\rm gr}$
and the proved bounds (ACM27), (ACM35) show that the interval
contains $\mathcal B_k$. Comparing the arguments of the lower
maximum and upper minimum, using each inequality in (ACM35),
proves the complete order in the last line. All later rows
included gives $L_k^{\rm corr}=U_k^{\rm corr}=\widehat{\mathcal B}_k$,
so both endpoints become exactly (HC9g). For any specified
subsets their complete width obeys
\[
\begin{aligned}
 0\leq b_{p,s}^{\rm corr,+}-b_{p,s}^{\rm corr,-}
 &\leq\min\{2e_{p,\Gamma},\,
 U_k^{\rm corr}(J)-L_k^{\rm corr}(J)
                     +g_{s,\rm gr}^+-g_{s,\rm gr}^-\}\\
 &\leq\min\{2e_{p,\Gamma},\,
 U_k^{\rm corr}(J)-L_k^{\rm corr}(J)+2e_{s,\rm gr}\}.
\end{aligned}\tag{HC10rcw}
\]
The Gamma interval lies in its radius-$e_{p,\Gamma}$ signed
interval, and the graph correction lies in its radius-$e_{s,\rm gr}$
signed interval by their actual finite remainder proofs. Taking
an intersection cannot increase either width. Subtracting the
displayed graph lower endpoint from its upper endpoint gives
the middle expression, proving every term of this bound.

The complete residual substitution at (HC9)--(HC10) is
\[
\begin{aligned}
 L_{k;p,s}^{\rm ar,corr}&=\max\{0,T_k,b_{p,s}^{\rm corr,-}\},\\
 U_{k;p,s}^{\rm ar,corr}&=\min\{\widehat{\mathcal U}_k,b_{p,s}^{\rm corr,+}\},\\
 L_{k;p,s}^{\rm ar,corr}-T_k
 &\leq\mathcal R_k\leq U_{k;p,s}^{\rm ar,corr}-T_k,\\
 L_{k;p,s}^{\rm ar,corr}-\mathcal B_k^\Gamma
 &\leq\Delta_k^\Gamma
                    \leq U_{k;p,s}^{\rm ar,corr}-\mathcal B_k^\Gamma.
\end{aligned}\tag{HC10rc}
\]
Indeed intersect (HC9rc) with the original inequalities
$\mathcal B_k\geq0$, $\mathcal B_k\geq T_k$ and
$\mathcal B_k\leq\widehat{\mathcal U}_k$ proved in (HC5),
(HC7), (HC8). The order-preserving translations already proved
in (HC10r) give the two displayed receiving intervals with their
exact inverse additions. Their widths do not exceed (HC10rcw).
The inclusion of intervals is now exact-graph inside correlation
inside relation-row inside Gamma; monotonicity of the maximum,
minimum and the two translations proves the same inclusion
at both residuals. All original lower-bound pieces and all
negative signed corrections remain present.


\subsection{The holonomy energy is tied to the same endpoint contractions}

Fix an actual holonomy $\theta$ and period $L$.  In this paragraph all
metrics, lifts and volumes carry that same phase.  Put
\[
 \begin{gathered}
 e_N=\log(V_N/V_{N+1}),\quad
 \beta_N=\lambda_{\max}
 (G_{N+1}^{-1/2}G_NG_{N+1}^{-1/2}),\\
 J_N^{D}=\|D_{L,\theta}R_NG_N^{-1/2}\|_{\rm HS}^{2},\quad
 \mathcal E_N=\|\mathscr B_NG_N^{-1/2}\|_{\rm HS}^{2}.
 \end{gathered}
\]
These use original source moments through degree $2N+2$.  They obey
\[
 \boxed{
 \frac{\mathcal L^2}{2}\leq\mathcal E_N
 \leq(1+\sqrt{\beta_N})^2J_N^D,
 \qquad
 e_N\geq2\log\max\left\{1,
 \frac{\mathcal L}{\sqrt{2J_N^D}}-1\right\}.}
 \tag{HC12}
\]
The denominator is positive: integration by parts on the stated
twisted domain gives
$\|D_{L,\theta}f\|^2=c^2\|f\|^2+\|\partial_rf\|^2$, whence
$J_N^D\geq c^2q>0$.

To prove the upper inequality, $D_{L,\theta}R_Nv$ is an original
degree-$(N+1)$ representative of $Av$.  The same quotient minimum
therefore gives
\[
 A^*G_{N+1}A\preceq R_N^*D_{L,\theta}^*D_{L,\theta}R_N.
\]
Use $G_N\preceq\beta_NG_{N+1}$, sandwich by $G_N^{-1/2}$, and take
traces.  This gives
$\|R_NAG_N^{-1/2}\|_{\rm HS}^2\leq\beta_NJ_N^D$.
The triangle inequality in $\mathscr B_N=D_{L,\theta}R_N-R_NA$
proves the upper inequality in (HC12); (HT14) proves the lower.
All eigenvalues in the definition of $\beta_N$ are at least one by
degree monotonicity.  In fact the original rank-one kernel update proves
the exact identity
\[
 e_N=\log\beta_N,\qquad
 \beta_N=1+\frac{b_{N+1}^*G_Nb_{N+1}}{\omega_{N+1}}.
 \tag{HC12a}
\]
Here $b_{N+1}$ and $\omega_{N+1}$ are the remainder column and monic
squared norm for this same original phase measure.  Indeed
$K_{N+1}=K_N+b_{N+1}b_{N+1}^*/\omega_{N+1}$, so
\[
 K_N^{-1/2}K_{N+1}K_N^{-1/2}=I+vv^*,\qquad
 v=K_N^{-1/2}b_{N+1}/\sqrt{\omega_{N+1}}.
\]
The relative matrix defining $\beta_N$ is
$K_{N+1}^{1/2}K_N^{-1}K_{N+1}^{1/2}=X^*X$ for the invertible
$X=K_N^{-1/2}K_{N+1}^{1/2}$.  The displayed matrix $I+vv^*$ is
$XX^*$, and $X(X^*X)X^{-1}=XX^*$ proves they have identical eigenvalues.
These are $1$ with multiplicity at least $q-1$ and $1+\|v\|^2$.
Taking the determinant gives
$V_N/V_{N+1}=\det K_{N+1}/\det K_N=1+\|v\|^2=\beta_N$,
which proves (HC12a), including a zero update if one occurs.
Rearranging the two energy bounds now proves the second inequality
of (HC12) with this exact endpoint identity.

The same weights as in (HC2) reconstruct the actual phase volume:
\[
 \boxed{
 \mathcal B_{k,\theta}=\sum_{N=q-1}^{2q-1}w_N e_N
 \geq2\sum_{N=q-1}^{2q-1}w_N
 \log\max\left\{1,\frac{\mathcal L}{\sqrt{2J_N^D}}-1\right\}.}
 \tag{HC13}
\]
Thus a small original derivative-source energy at a degree has the
explicit endpoint contraction cost (HC12); that cost enters the same
four-window budget, with its actual multiplicity.

Finally the supplied holonomy construction allows any fixed
$0<\eta<1$ by choosing the original period from its computed weighted
source comparison at degree $2q+2$.  It proves simultaneously
$(1-\eta)G_N\preceq G_{N,\theta}\preceq(1+\eta)G_N$ at all these
degrees.  For completeness, apply the comparison before quotient
minimization on the common coefficient space $H=\mathcal P_{2q}$.
The supplied weighted source estimate, written explicitly in (PSC.33),
is $(1-\eta)M_{2q}\preceq M_{2q,\theta}\preceq(1+\eta)M_{2q}$.
Set $H_\theta(t)=(1-t)M_{2q}+tM_{2q,\theta}$,
$C^\theta=H_\theta(t)^{-1}(M_{2q,\theta}-M_{2q})$, and define
$P_N^\theta,Q_N^\theta$ by (AT11) with $M_*=H_\theta(t)$ and
the unchanged inclusions $I_N,B_N$. Then
$U^\theta=Q_{2q-1}^\theta+Q_{2q}^\theta-Q_q^\theta$ and
$W^\theta=P_{2q-1}^\theta-P_{q-1}^\theta+P_{2q}^\theta-P_q^\theta$.
Substitution proves all projector identities and the exact signed
trace formula (AT15) on this distinct interpolation. Both positive
forms have the same original remainder fibres. Thus (AT18a)--(AT21)
defines its own control $\mathfrak C_\theta$ and proves
\[
\begin{aligned}
 |\mathcal B_{k,\theta}-\mathcal B_k|
 &\le J_\theta\le\mathfrak C_\theta
 :=\int_0^1\min\{S_{\rm spec}^\theta(t),
        \tfrac12d_\theta(t)\|U^\theta-W^\theta\|_{1,H_\theta(t)}\}\,dt\\
 &\le(2q-1)\log\frac{1+\eta}{1-\eta}.
\end{aligned}\tag{HC14}
\]
Indeed every eigenvalue of $M_{2q}^{-1}M_{2q,\theta}$ lies in
$[1-\eta,1+\eta]$ by the source sandwich, and the full spectra
and trace norms in the finite AT proof are on the displayed
polynomial flags. This is an explicit source comparison; taking
determinants only after quotient minimization gives the preceding
$2q$ allowance as a coarser consequence. The interpolated forms
are not identified with either Gamma reference.
On these same phase-source operators define
$S_{\rm HS}^\theta$ by the exact centered trace formula (AT21a),
and $J_\theta=\int_0^1\min\{S_{\rm spec}^\theta,
S_{\rm tr}^\theta,S_{\rm HS}^\theta\}\,dt$.
The same finite Hermitian Cauchy--Schwarz proof gives the first
entry in (HC14); hence $J_\theta\le\mathfrak C_\theta$.
The positive atomic phase measure also meets the moment hypothesis
proved above. Define $I_\theta$ by (AT21a) on the exact path
$H_\theta(t)$; its original $B_\theta(t)$ is strictly positive by
the same nonzero moment $\int|\chi|^2d\mu_t$. Put
\[
\begin{aligned}
\mathfrak U_{\rm ar}&=\min\{\widehat{\mathcal U}_k,\mathfrak b_*^+\},\\
\mathfrak U_\theta&=\min\{\mathfrak U_{\rm ar}+J_\theta,
 \Psi_a(H_a(\mathfrak U_{\rm ar})+I_\theta)\},\qquad a=2q-1.
\end{aligned}\tag{HC14a}
\]
Then (HC13)'s right side is at most $\mathfrak U_\theta$.
Indeed $0<\mathcal B_k\le\mathfrak U_{\rm ar}$ by the original
upper estimates. Apply (AT21e) on $H_\theta(t)$, then use
monotonicity of $\Psi_a\circ(H_a+I_\theta)$ in its positive
initial value. This proves both upper entries and their minimum.
This substitution uses (AT19) on the original arithmetic form
before the explicit source-phase trace calculation; it does not identify
the sampled phase metric with the original one.

\paragraph{The original source-phase map and its signed refinement.}
We now refine the same bound (HC14a), retaining its original
$J_\theta,I_\theta,\mathfrak C_\theta$ and both upper entries.
Let $I_{2q}^{2q+2}:\mathcal P_{2q}\hookrightarrow\mathcal P_{2q+2}$
be literal coefficient inclusion. The two original forms on
$H=\mathcal P_{2q}$ are
\[
\begin{aligned}
 M_0^\theta
 &=(I_{2q}^{2q+2})^*M_{2q+2}I_{2q}^{2q+2}
   =\int_{\mathbb R}v(u)^*v(u)m_{h,k}(u)\,du,\\
 M_1^\theta
 &=(I_{2q}^{2q+2})^*M_{2q+2,\theta}I_{2q}^{2q+2}
   =\frac{2\pi}{L}\sum_{r\in\mathbb Z}
                 m_{h,k}(u_r)v(u_r)^*v(u_r),\\
 v(u)&=(1,c+iu,\ldots,(c+iu)^{2q}),\qquad
 c=k/2,\qquad u_r=(2\pi r+\theta)/L .
\end{aligned}\tag{HC14b}
\]
By (HT3)--(HT7), these are exactly the squared norms of
$\mathcal U_k\mathcal V|_H$ and
$\mathcal Z_{L,\theta}\mathcal U_k\mathcal V|_H$.
Here $\mathcal U_k$ is the half-density map (HT4), whose source
and target include every relative variable; the earlier scalar
upper bound with this notation is still the expression in (HC8).
The original relative variables, full Taylor unit, coordinate,
phase, and factor $2\pi/L$ are present in these two source maps.
To specify the period cost without colliding with $a=2q-1$,
write $a_{\rm wt}>0$ for the weighted-source exponent of (PSC.32).
Its actual degree-$2q+2$ cost is
\[
 \kappa_{a_{\rm wt},2q+2}
 =\sup_{v\ne0}
 \frac{v^*(M_{a_{\rm wt},+,2q+2}+M_{a_{\rm wt},-,2q+2})v}
      {v^*M_{2q+2}v}.
\]
The weighted Grams retain the squared weights $x_k^{\pm2a_{\rm wt}}$.
Retain the originally fixed period $L$ whenever
$\kappa_{a_{\rm wt},2q+2}/(e^{a_{\rm wt}L}-1)\le\eta$.
One exact admissible choice is
$L=a_{\rm wt}^{-1}\log(1+\kappa_{a_{\rm wt},2q+2}/\eta)$;
the comparison below holds for every retained period satisfying
the displayed inequality.
Equation (PSC.33), restricted by $I_{2q}^{2q+2}$, proves
$(1-\eta)M_0^\theta\preceq M_1^\theta\preceq(1+\eta)M_0^\theta$.
For $q-1\le N\le2q$ the remainder $J_N$ is surjective; taking
the two infima over its identical nonempty fibre $J_NP=v$
also proves the quotient sandwiches used above.

Put $M_x^\theta=(1-x)M_0^\theta+xM_1^\theta$ for $0\le x\le1$.
At $x=t$ this is precisely $H_\theta(t)$ in (HC14).
Use $I_N:\mathcal P_N\to H$ and the unpadded original relation
$B_N^{\rm AW}:\mathcal P_{N-q}\to\mathcal P_N$, $P\mapsto\chi P$.
Its included map is $B_N^\theta=I_NB_N^{\rm AW}$, and the
exact projections and operators are
\[
\begin{aligned}
 P_N^\theta&=I_N(I_N^*M_x^\theta I_N)^{-1}I_N^*M_x^\theta,\\
 Q_N^\theta&=B_N^\theta((B_N^\theta)^*M_x^\theta B_N^\theta)^{-1}
                                  (B_N^\theta)^*M_x^\theta,\\
 U_x^\theta&=Q_{2q-1}^\theta+Q_{2q}^\theta-Q_q^\theta,\\
 W_x^\theta&=P_{2q-1}^\theta-P_{q-1}^\theta+P_{2q}^\theta-P_q^\theta,\\
 C_x^\theta&=(M_x^\theta)^{-1}(M_1^\theta-M_0^\theta),\qquad
 D_\theta=(M_0^\theta)^{-1}M_1^\theta .
\end{aligned}\tag{HC14c}
\]
At $N=q-1$ the relation domain is zero and the second formula
is the unique zero map $Q_{q-1}^\theta=0$.
These are exactly (ACM4)--(ACM5) after the stated substitution,
and exactly the projections already used in (HC14).
The source and relation ranges are the same nested polynomial
flags, so both $U_x^\theta,W_x^\theta$ have full spectrum
$0^{[q]},1^{[2]},2^{[q-1]}$. Differentiating the actual
canonical section as in (ISM18) and taking the four determinant
signs gives
\[
 F_\theta'(x)=\operatorname{Tr}(C_x^\theta(U_x^\theta-W_x^\theta)),
 \quad F_\theta(0)=\mathcal B_k,\quad
 F_\theta(1)=\mathcal B_{k,\theta}.
\]
The polynomial $\chi$ has positive norm for the continuous
measure and for the infinite positive lattice measure (HT6).
Their convex mixture has the same property. The original
$\chi^{\#_k}$ proof (AT21a)--(AT21c) therefore gives
$F_\theta(x)>0$ on this closed path. In particular every
nonlinear denominator is bounded away from zero for this fixed
source comparison.

Apply the pointwise dictionary preceding (HC9a) on this path.
The additional trace inequalities
$S_{\rm tr}^\theta\le S_{\rm ov}^\theta$ and
$S_{\rm tr}^\theta\le S_{\rm sin}^\theta\le S_{\rm ang}^\theta$
prove that the complete (ACM11a) minima are precisely the
$J_\theta,I_\theta$ already defined above. In particular both
complete angle lists, the overlap majorant and the centered
Hilbert--Schmidt expression are retained in these identities.
Let $j_{p,\theta},e_{p,\theta}$ be the (ACM15) construction with
\[
 (M_x,C_x,U_x,W_x,D,F(x))
 \longmapsto
 (M_x^\theta,C_x^\theta,U_x^\theta,W_x^\theta,D_\theta,F_\theta(x)).
\]
The integer $p\ge0$ is its finite-series degree, while the
original period is still $L$. If
$b_1^\theta\le\cdots\le b_n^\theta$ are every eigenvalue of
$D_\theta$ in $M_0^\theta$, $n=2q+1$, write
\[
 \vartheta_\theta=\frac{b_n^\theta-b_1^\theta}{b_n^\theta+b_1^\theta},
 \qquad
 K_{D_\theta}=\sum_{j=1}^n
       \left(\frac{|b_j^\theta-1|}{\min\{1,b_j^\theta\}}\right)^2 .
\]
The finite identity
$C_x^\theta-(C_x^\theta)^{[p]}
 =(H_x^{\circ,\theta})^{p+1}C_x^\theta$
is the identical (ACM15) geometric sum on these specified
operators. Subtracting the trace mean of this residual leaves
its pairing with $U_x^\theta-W_x^\theta$ unchanged.
All its factors are real functions of $D_\theta$, hence
self-adjoint in $M_x^\theta$. The Hilbert--Schmidt inequality
therefore gives
\[
\begin{gathered}
 j_{p,\theta}-e_{p,\theta}
 \le\mathcal B_{k,\theta}-\mathcal B_k
 \le j_{p,\theta}+e_{p,\theta},\\
 0\le e_{p,\theta}\le
 \vartheta_\theta^{\,p+1}\sqrt{K_{D_\theta}(8q-6)}
 \xrightarrow[p\to\infty]{}0 .
\end{gathered}\tag{HC14d}
\]
For completeness the first factor bounds every residual
eigenvalue because
$\|H_x^{\circ,\theta}\|_{M_x^\theta}\le\vartheta_\theta<1$ and
$|c_j^\theta(x)|\le|b_j^\theta-1|/\min\{1,b_j^\theta\}$.
Centering subtracts the nonnegative square of the trace mean.
Both positive operators dominate their range projections, whose
ranges intersect in dimension at least one. Thus
$\operatorname{Tr}(U_x^\theta W_x^\theta)\ge1$, and the full
spectra give
$\operatorname{Tr}((U_x^\theta-W_x^\theta)^2)
 =8q-4-2\operatorname{Tr}(U_x^\theta W_x^\theta)\le8q-6$.
This proves the displayed constant on the original metric.

The source sandwich places every $b_j^\theta$ in
$[1-\eta,1+\eta]$. It consequently gives
$\vartheta_\theta\le\eta$ and
$K_{D_\theta}\le(2q+1)\eta^2/(1-\eta)^2$: the first follows
by increasing the larger endpoint and decreasing the smaller
one in their positive ratio, and each summand of the second
has numerator at most $\eta^2$ and denominator at least
$(1-\eta)^2$. Hence the explicit further bound is
\[
 e_{p,\theta}\le
 \frac{\eta^{p+2}}{1-\eta}\sqrt{(2q+1)(8q-6)} .
\]
This holds for every phase at each retained period satisfying
the displayed original degree-$2q+2$ source-cost inequality, retaining
the complete $q=[1+k(m-1)](k+1)^2$ dependence.
Retaining their entire ordered list, the
same signed derivative and the exact derivative of $H_a$ give
\[
\begin{aligned}
 |\mathcal B_{k,\theta}-\mathcal B_k|
 &\le J_\theta\le
 \min\{\mathfrak C_\theta,\mathcal L_q(D_\theta)\}
 \le(2q-1)\log\frac{1+\eta}{1-\eta},\\
 |H_{2q-1}(\mathcal B_{k,\theta})-H_{2q-1}(\mathcal B_k)|
 &\le I_\theta\le\log\frac{1+\eta}{1-\eta},\\
 \mathcal L_q(D_\theta)
 &=\log\frac{b_{q+2}^\theta}{b_q^\theta}
   +2\sum_{j=1}^{q-1}\log\frac{b_{2q+2-j}^\theta}{b_j^\theta}.
\end{aligned}\tag{HC14e}
\]
Indeed integration of each ordered
$c_j^\theta(x)=(b_j^\theta-1)/(1-x+xb_j^\theta)$ gives
$\log b_j^\theta$. The total coefficient of the resulting
logarithmic ratios is $1+2(q-1)=2q-1$.
Dividing the original pointwise minimum by
$(2q-1)\sqrt{1-e^{-F_\theta(x)/(2q-1)}}$
and integrating proves the second line. These are the same
additive and nonlinear bounds as (HC14)--(HC14a), now
accompanied by the constructed signed center (HC14d).

For any integers $p,r\ge0$, define
\[
\begin{aligned}
 U_{k,p}^{\rm ar}
 &=\min\{\widehat{\mathcal U}_k,\beta_{p,\Gamma}^+\},\\
 U_{k,\theta;p,r}
 &=\min\bigl\{
 U_{k,p}^{\rm ar}+J_\theta,\,
 \Psi_{2q-1}(H_{2q-1}(U_{k,p}^{\rm ar})+I_\theta),\\
 &\hspace{32mm}U_{k,p}^{\rm ar}+j_{r,\theta}+e_{r,\theta}
 \bigr\}.
\end{aligned}\tag{HC14f}
\]
The actual $\mathcal B_k$ is positive and at most
$U_{k,p}^{\rm ar}$ by (HC9a)--(HC10a).
The additive and signed expressions increase with this initial
scalar. For the fixed computed $I_\theta$, both $H_{2q-1}$
and $\Psi_{2q-1}$ are increasing on the displayed domains.
Applying the three proved phase upper bounds and substituting
this initial upper value proves
\[
 2\sum_{N=q-1}^{2q-1}w_N
 \log\max\left\{1,\frac{\mathcal L}{\sqrt{2J_N^D}}-1\right\}
 \le\mathcal B_{k,\theta}
 \le U_{k,\theta;p,r}\le\mathfrak U_\theta .
\tag{HC13a}
\]
For the final inequality (HC9a) gives
$U_{k,p}^{\rm ar}\le\mathfrak U_{\rm ar}$, and monotonicity
of the first two upper expressions bounds their minimum by
the original (HC14a). The third entry can only lower it.
This proof holds the actual phase-path integrals fixed during
the scalar comparison; it does not replace their source path.
The weights $w_N$, derivative-source energies $J_N^D$, and every
boundary term in (HC12)--(HC13) remain in their original phase
metric. No fixed-source residual limit is promoted to a
uniform estimate in the growing tensor order.

The source comparison also retains its exact cochain map.
For two minimum sections $\rho_0,\rho_1$ on this common source,
$(\rho_1-\rho_0)v=\chi Q_v$ by monic division. Successive
division in $s_1,\ldots,s_k$ gives the original polynomials
$Q_{v,i}$ with
$\chi(\sum s_i)Q_v(\sum s_i)=\sum_i h(s_i)Q_{v,i}(\mathbf s)$.
With the original $h(D)F_h=\Theta\phi_*$, the actual primitive
is
\[
 K_v=\sum_{i=1}^k(-1)^{i-1}Q_{v,i}(D_1,\ldots,D_k)
 (F_h^{\otimes(i-1)}\otimes\phi_*
                         \otimes F_h^{\otimes(k-i)}).
\]
The tensor differential contributes the same $(-1)^{i-1}$,
so $dK_v=\mathcal V((\rho_1-\rho_0)v)$, as in (ISM38)--(ISM41).
For the two-leg complex with differential
$(\phi,\widehat\psi)\mapsto\Theta(\phi-\mathcal F^{-1}\widehat\psi)$,
the degree-zero map $j_+(\phi)=(\phi,0)$ and degree-one identity
form a cochain map because their differential is exactly
$\Theta\phi$. Tensor these specified maps and use the original
coefficient insertions into the fixed full carrier of (ISM42).
This defines $\widetilde K_v$ and preserves its displayed top
differential. At a declared support $\lambda=(W,S;A)$, apply
these maps coordinatewise. With the original coefficient
injections $\iota_i$, the exact top maps and their common
inverse-image domain are
\[
\begin{aligned}
 F_{a,A}:E^{\oplus A}&\longrightarrow\mathsf C_{\rm full}^k,
 &F_{a,A}((v_i)_{i\in A})&=\sum_{i\in A}\iota_i\mathcal V\rho_av_i
 \quad(a=0,1),\\
 E_\lambda&=
 F_{0,A}^{-1}(\mathsf C_\lambda^k)
 \cap F_{1,A}^{-1}(\mathsf C_\lambda^k).
\end{aligned}
\]
This is a linear subspace. For $\mathbf v=(v_i)_{i\in A}\in E_\lambda$,
put $\widetilde K_{\mathbf v}=\sum_{i\in A}\iota_i j_+^{\otimes k}K_{v_i}$,
where the tensor cochain map has the degree-one identities already
specified. Its top differential is $F_{1,A}(\mathbf v)-F_{0,A}(\mathbf v)$,
which lies in $\mathsf C_\lambda^k$ and is a boundary in the full
carrier. The exact residual map on this specified domain is
\[
 E_\lambda\longrightarrow
 \frac{d\mathsf C_{\rm full}^{k-1}\cap\mathsf C_\lambda^k}
      {d\mathsf C_\lambda^{k-1}},\qquad
 \mathbf v\longmapsto[d\widetilde K_{\mathbf v}]_\lambda .
\]
Its inclusion-induced map to the kernel of
$\mathsf C_\lambda^k/d\mathsf C_\lambda^{k-1}
 \to\mathsf C_{\rm full}^k/d\mathsf C_{\rm full}^{k-1}$
and its inverse both use the same representative, proving the
exact isomorphism (ISM42). The full two-leg degree-zero
diagonal and the source images are those of (ISM41)--(ISM43).
For independent coefficient masks the linear maps act in each
active coordinate; zero insertion commutes because every new
coordinate maps $0$ to $0$ with its original outer label.
This calculation retains the proper-source residual and the
empty mask's outer label throughout.
\paragraph{The graph refinement in this original phase-energy budget.}
Retain the domain $\gamma>2$ of (HC9g), the actual original period
and phase above, and every previously calculated $J_\theta,I_\theta,
j_{r,\theta},e_{r,\theta}$. For $p,s,r\geq0$ put
\[
 \begin{aligned}
 U_{k;p,s}^{\rm ar,gr}
 &=\min\{\widehat{\mathcal U}_k,b_{p,s}^{\rm gr,+}\},\\
 U_{k,\theta;p,s,r}^{\rm gr}
 &=\min\bigl\{U_{k;p,s}^{\rm ar,gr}+J_\theta,\\
 &\qquad\Psi_{2q-1}\bigl(
       \mathscr H_{2q-1}(U_{k;p,s}^{\rm ar,gr})+I_\theta\bigr),\\
 &\qquad U_{k;p,s}^{\rm ar,gr}+j_{r,\theta}+e_{r,\theta}\bigr\}.
 \end{aligned}
 \tag{HC14g}
\]
The original arithmetic moment proof at (ACM11a) gives strict
positivity of $\mathcal B_k$; (HC10g) and (HC9g) then give
$0<\mathcal B_k\leq U_{k;p,s}^{\rm ar,gr}\leq U_{k,p}^{\rm ar}$.
For each of the three fixed phase controls the upper expression
is increasing in its initial scalar: the first and third are
affine of slope one, and the second is the composite of the
two increasing functions already differentiated in (ACM14).
Applying their actual phase inequalities and this scalar
comparison proves the propagated bound at (HC13a):
\[
 \begin{aligned}
 2\sum_{N=q-1}^{2q-1}w_N
 \log\max\left\{1,\frac{\mathcal L}{\sqrt{2J_N^D}}-1\right\}
 &\leq\mathcal B_{k,\theta}\\
 &\leq U_{k,\theta;p,s,r}^{\rm gr}
 \leq U_{k,\theta;p,r}.
 \end{aligned}
 \tag{HC13g}
\]
The original $w_N$, $J_N^D$, phase source, twisted derivative
domain and boundary residual in (HC12)--(HC13) are unchanged.
The graph data enter through the proved signed bound for the
initial arithmetic volume. The inequality does not identify
the two-component graph source with the phase source.

\paragraph{The relation-row interval propagated into the phase budget.}
Keep the original phase, circumference, source-path matrices and
the exact integrals $J_\theta,I_\theta,j_{r,\theta},e_{r,\theta}$
from (HC14b)--(HC14g). Use $U_{k;p,s}^{\rm ar,rel}$ defined in
(HC10r), and define the complete three-entry upper value
\[
\begin{aligned}
 U_{k,\theta;p,s,r}^{\rm rel}
 =\min\bigl\{&U_{k;p,s}^{\rm ar,rel}+J_\theta,\\
 &\Psi_{2q-1}\bigl(\mathscr H_{2q-1}
            (U_{k;p,s}^{\rm ar,rel})+I_\theta\bigr),\\
 &U_{k;p,s}^{\rm ar,rel}+j_{r,\theta}+e_{r,\theta}\bigr\}.
\end{aligned}\tag{HC14r}
\]
By (HC9r)--(HC10r) and the strict original-moment positivity,
\[
 0<\mathcal B_k\leq U_{k;p,s}^{\rm ar,gr}
             \leq U_{k;p,s}^{\rm ar,rel}\leq U_{k,p}^{\rm ar}.
\]
For this fixed phase path the first and third functions of an
initial scalar $x$ in (HC14r) are affine of slope one. The second
is the composite of increasing $\mathscr H_{2q-1}$ on $x>0$
and increasing $\Psi_{2q-1}$ on its nonnegative argument; their
derivatives and inverse domains were computed in (ACM14).
The actual phase inequalities (HC14d)--(HC14e) use
$x=\mathcal B_k$. Substituting the larger scalar upper endpoints
and applying each monotonicity proves
\[
\begin{aligned}
2\sum_{N=q-1}^{2q-1}w_N
 \log\max\left\{1,\frac{\mathcal L}{\sqrt{2J_N^D}}-1\right\}
 &\leq\mathcal B_{k,\theta}\leq U_{k,\theta;p,s,r}^{\rm gr}\\
 &\leq U_{k,\theta;p,s,r}^{\rm rel}\leq U_{k,\theta;p,r}.
\end{aligned}\tag{HC13r}
\]
The left inequality is the unchanged original phase-energy
identity (HC13), with the complete multiplicities, twisted
derivative domain and boundary terms. The comparison holds the
actual path controls fixed; the scalar substitution does not
change their integration path or measure. It also retains the
entire cochain map and proper-source residual following (HC13a):
the only new maps here are the two real translations in (HC10r)
and these three functions of their proved arithmetic upper
endpoint. The exact-graph phase bound (HC13g) remains the smaller
one at the same $p,s,r$, as the displayed inequalities prove.


\paragraph{The correlation interval in the three original phase functions.}
With the same fixed original phase controls as (HC14r), set
\[
\begin{aligned}
 U_{k,\theta;p,s,r}^{\rm corr}
 =\min\bigl\{&U_{k;p,s}^{\rm ar,corr}+J_\theta,\\
 &\Psi_{2q-1}\bigl(\mathscr H_{2q-1}
          (U_{k;p,s}^{\rm ar,corr})+I_\theta\bigr),\\
 &U_{k;p,s}^{\rm ar,corr}+j_{r,\theta}+e_{r,\theta}\bigr\}.
\end{aligned}\tag{HC14rc}
\]
Equation (HC9rc) and the arithmetic minimum in (HC10rc) give
$0<\mathcal B_k\leq U_{k;p,s}^{\rm ar,gr}
\leq U_{k;p,s}^{\rm ar,corr}\leq U_{k;p,s}^{\rm ar,rel}$.
The three increasing functions proved in (HC14r) therefore give
\[
\begin{aligned}
 2\sum_{N=q-1}^{2q-1}w_N
 \log\max\left\{1,\frac{\mathcal L}{\sqrt{2J_N^D}}-1\right\}
 &\leq\mathcal B_{k,\theta}\leq U_{k,\theta;p,s,r}^{\rm gr}\\
 &\leq U_{k,\theta;p,s,r}^{\rm corr}
 \leq U_{k,\theta;p,s,r}^{\rm rel}\leq U_{k,\theta;p,r}.
\end{aligned}\tag{HC13rc}
\]
This proves the full propagation into the original phase-energy
budget by applying each actual phase inequality before taking
their minimum. The source-path integrals, circumference condition,
twisted derivative domain, complete weights, boundary terms and
proper-source cochain residual are those already proved above.
Including all later rows makes the correlation upper bound equal
to the exact-graph bound (HC14g) at the same $p,s,r$. No tensor
uniform rate is inferred from this finite exactness statement.


For the averaged original phase metrics the already proved stronger
comparison $(1-\eta^2)G_N\preceq\overline G_N\preceq G_N$ gives instead
$|\overline{\mathcal B}_k-\mathcal B_k|
\leq2q\log(1/(1-\eta^2))$.  These changes are of order $q$ for a fixed
$\eta$, whereas the compatible width (HC11) has positive quadratic
order.  The exact Laplacian residual (HT17) remains present at every
phase.  No displayed upper estimate for its norm, or for the original
derivative moments in (HC12), contradicts these lower bounds.

\paragraph{The graph-to-average cost and its complete mixed remainder.}
On the same graph-source domain $\gamma>2$, fix one original
circumference for all four endpoints. The inverse-form statements
below use the already proved degree-$2q+2$ original source
comparison above. For the graph map $L_{\rm obs}$ of (ACM18),
write $L_{\rm circ}$ for this circumference to avoid a collision.
At each $N\in(q-1,q,2q-1,2q)$ use the original arithmetic
minimum section $C_N$ and the relation map $B_N$ of PGB.3.
Put $H_N=B_N^*M_NB_N$, $T_N^{\rm def}=D_NB_N$,
$H_{+,N}=H_N+(T_N^{\rm def})^*T_N^{\rm def}$, and retain
\[
 Q_{D,N}=I-T_N^{\rm def}H_{+,N}^{-1}(T_N^{\rm def})^*,
 \quad
 \mathcal E_N^{\rm aug}=C_N^*D_N^*Q_{D,N}D_NC_N,
 \quad
 \mathscr D_N^{\rm hol}=G_N-\overline G_N.
 \tag{HC16}
\]
For zero relation dimension the maps through that domain are
zero and $Q_{D,N}=I$. PGB.3--PGB.10 proves, by the displayed
completed squares and the original phase section difference,
\[
 \begin{aligned}
 \widehat G_N&=G_N+\mathcal E_N^{\rm aug},\qquad
 \mathcal E_N^{\rm aug}>0,\\
 \mathscr D_N^{\rm hol}
 &=\int(C_N-C_{N,\theta})^*M_{N,\theta}
                     (C_N-C_{N,\theta})\frac{d\theta}{2\pi}
 \succeq0,\\
 \widehat G_N-\overline G_N&=
      \mathcal E_N^{\rm aug}+\mathscr D_N^{\rm hol}.
 \end{aligned}
 \tag{HC17}
\]
The exact map whose Gram is the last difference is
\[
 x\longmapsto\left(Q_{D,N}^{1/2}D_NC_Nx,
           [\theta\mapsto(C_N-C_{N,\theta})x]\right)
 \in\mathscr H_k\oplus
 L^2(\Theta;\operatorname{im}B_N,M_{N,\theta}d\theta/(2\pi)).
 \tag{HC18}
\]
Its squared norm is (HC17), and its kernel is zero by the
injectivity of $D_NC_N$ proved from the full leading coefficient
in PGB.6. The cochain map inducing the identity contraction
$(E,\widehat G_N)\to(E,\overline G_N)$ sends each polynomial
and each relation to its constant phase field. It commutes
with the inclusion differential and preserves $J_N$; its
norm identity is $\int p^*M_{N,\theta}p\,d\theta/(2\pi)
=p^*M_Np\leq p^*\widehat M_Np$.
Thus the finite comparison is accompanied by its original
source map and its complete nonzero metric loss.

To carry this calculation into the four-volume use site define
the exact individual determinant costs
\[
 \begin{aligned}
 a_N^{\rm aug}&=\log\frac{\det\widehat G_N}{\det G_N}
 =\log\det(I+G_N^{-1/2}\mathcal E_N^{\rm aug}G_N^{-1/2})>0,\\
 h_N^{\rm hol}&=\log\frac{\det G_N}{\det\overline G_N}
 =\log\det(I+\overline G_N^{-1/2}
                  \mathscr D_N^{\rm hol}\overline G_N^{-1/2})
 \geq0,\\
 a_N^{\rm av}&=\log\frac{\det\widehat G_N}{\det\overline G_N}
                    =a_N^{\rm aug}+h_N^{\rm hol}.
 \end{aligned}
 \tag{HC19}
\]
These equalities follow by determinant multiplication; no
commutation of their positive matrices is required. List the
four endpoints as $(N_1,N_2,N_3,N_4)=(q-1,q,2q-1,2q)$ with
$\epsilon=(1,1,-1,-1)$, retaining repeated entries if they
occur. The exact signed identities are
\[
 \begin{aligned}
 \Gamma_k^{\rm gr}&=\sum_{j=1}^4\epsilon_j a_{N_j}^{\rm aug},\\
 H_k^{\rm hol}&:=\sum_{j=1}^4\epsilon_j h_{N_j}^{\rm hol}
                       =\mathcal B_k-\overline{\mathcal B}_k,\\
 \widehat{\mathcal B}_k-\overline{\mathcal B}_k
 &=\Gamma_k^{\rm gr}+H_k^{\rm hol}
   =\sum_{j=1}^4\epsilon_j a_{N_j}^{\rm av}.
 \end{aligned}
 \tag{HC20}
\]
For the original source sandwich chosen above,
$(1-\eta^2)G_N\preceq\overline G_N\preceq G_N$ implies
$0\leq h_N^{\rm hol}\leq q\log(1/(1-\eta^2))$.
There are two positive and two negative signs, so
\[
 |H_k^{\rm hol}|\leq2q\log\frac1{1-\eta^2},\qquad
 \left|\widehat{\mathcal B}_k-\overline{\mathcal B}_k
                       -\Gamma_k^{\rm gr}\right|
 \leq2q\log\frac1{1-\eta^2}.
 \tag{HC21}
\]
This proves exactly where the earlier order-$q$ allowance
enters. The augmentation's signed term is retained through
(ACM23)--(ACM27), rather than assigned that allowance.
The separate positivity of the four costs in (HC19) does not
give a sign to their complete signed sum.

The graph volume appearing in (HC9g) also has its full
mixed-source expression. Retain $\Phi(1/2+it)=A+iq_\Phi(t)$,
the density $r_k$, function $g_x=z(u)x/r_k(u)$, and maps
$Y_\perp,\mathscr J_k,\widehat R_N$ exactly as in PGB.14--PGB.22.
In particular $\Omega_k=Y_\perp^*Y_\perp>0$ and
$f_Nx=\iota\widehat R_Nx-g_x$, with $\iota P(u)=P(k/2+iu)$.
The complete squared-norm identity PGMB.5 is
\[
 \begin{aligned}
 \widehat G_N&=K_{\rm low}
             +(\mathscr D_k^{\rm mix})^*\mathscr D_k^{\rm mix}
             +f_N^*f_N=\Omega_k+f_N^*f_N,\\
 (\mathscr D_k^{\rm mix}x)(\mathbf t)
 &=\frac{\overline{\tau_k}F_x(\mathbf s)}{1+|\tau_k|^2}-g_x(u).
 \end{aligned}
 \tag{HC22}
\]
The target norm of the second line is the full weighted space
$L^2((1+|\tau_k|^2)\prod_iw_h(t_i)dt_i)$.
Thus it retains $1+k^2A^2+(\sum_iq_\Phi(t_i))^2$, the conjugate
phase and every tensor variable. PGMB.2 evaluates its kernel:
it is zero for $k\geq2$, while the map vanishes for $k=1$.
Its first term has the original realization
$L_{\rm low}x=(\prod_i v_h(s_i))F_x/\sqrt{1+|\tau_k|^2}$,
with Gram $K_{\rm low}$. The map
$x\mapsto(L_{\rm low}x,\mathscr D_k^{\rm mix}x,f_Nx)$
has precisely the Gram $\widehat G_N$ in (HC22), proving
the source realization of every term.

Define its finite positive tail costs by
\[
 t_N=\log\det(I+\Omega_k^{-1/2}f_N^*f_N\Omega_k^{-1/2}).
 \tag{HC23}
\]
Equation (HC22) gives
$\log\det\widehat G_N=\log\det\Omega_k+t_N$.
The minimum over the nested original relation spaces proves
$\widehat G_{N+1}\preceq\widehat G_N$, hence
$t_{N+1}\leq t_N$; determinant monotonicity follows by
congruence with the smaller positive form and its eigenvalues.
PGB.22's affine approximation proves $f_N^*f_N\to0$ at fixed
$h,k$, so $t_N\downarrow0$ in that limit. The actual four
endpoints therefore give
\[
 \begin{aligned}
 \widehat{\mathcal B}_k
   &=(t_{q-1}-t_{2q-1})+(t_q-t_{2q}),\\
 0&\leq\widehat{\mathcal B}_k\leq t_{q-1}+t_q,\\
 \mathcal B_k&=t_{q-1}+t_q-t_{2q-1}-t_{2q}
                                         -\Gamma_k^{\rm gr}.
 \end{aligned}
 \tag{HC24}
\]
Substitution into (HC9g), (HC10g) and then (HC14g) is the
explicit original-source tail calculation for their graph
entry. In particular it retains all four tails and the
signed augmentation correction. Although the common
$\log\det\Omega_k$ cancels with $1+1-1-1=0$, that full
matrix remains in all four congruences (HC23). Its strictly
positive mixed contribution at every $k\geq2$ has not been
removed. The fixed-$k$ degree limit of $t_N$ does not supply
a rate for these four growing endpoints as $k$ varies.

Finally the completed receiving maps remain explicit. PGB.20--PGB.21
gives $\Xi(f,x)=\mathscr J_k f-Y_\perp x$ and
$j_\infty\Xi(f,x)=x$, where
$j_\infty=-\Omega_k^{-1}Y_\perp^*$ and
$\ker j_\infty=\mathscr J_kL^2(r_kdu)$.
The first-coordinate projection to the original unweighted
completed source is
$\pi_1\Xi(f,x)=\mathscr J_k^0 I_{r,m}(f+g_x)$.
It maps the closed graph relation into the original closed
unweighted relation, inducing $(E,\Omega_k)\to0$ with full
kernel $E$, exactly as proved in PGB.21b. On the graph quotient,
PGMB.7--PGMB.8 instead gives
$\mathfrak C_k\Xi(f,x)=(L_{\rm low}x,\mathscr D_k^{\rm mix}x)$
with kernel $\mathscr J_kL^2(r_kdu)$, and, for $k\geq2$,
\[
 (\Omega_k-K_{\rm low})^{-1}
          (\mathscr D_k^{\rm mix})^*\operatorname{pr}_2
                        \mathfrak C_k=j_\infty.
 \tag{HC25}
\]
The equality follows by the exact Gram in (HC22).
These are the maps from the same completed source to its
unweighted and mixed targets, with their complete kernels.
On each fixed coefficient mask they act coordinatewise;
represented zero and an empty present tuple keep their
original label, and external absence keeps its singleton
inverse image. The finite graph-to-average contraction
(HC18) and these completed maps use their stated different
domains. Every proper-source residual retained above
continues to use its original inverse-image domain and
boundary quotient; none of these Hilbert quotient calculations
declares that residual zero.


\subsection{The sharper finite comparison retains its exact error}

The original lower-envelope measure is
\[
 d\nu_k(u)=c_h\vartheta_h^{k-3}e^{-\alpha(k-3)}
 e^{-\alpha|u|}(k-2+|u|)^{-B}\,du.
\]
Let $V_N^\nu$ be its quotient volume for the identical full remainder
map.  The actual finite upper expression for the four-window volume is
\[
 \mathcal U_{\nu,k}=\log(V_{q-1}V_q)
                  -\log(V_{2q-1}^\nu V_{2q}^\nu),
\]
with the exact nonnegative error
\[
 \boxed{
 \mathcal U_{\nu,k}-\mathcal B_k
 =\log\frac{V_{2q-1}}{V_{2q-1}^\nu}
  +\log\frac{V_{2q}}{V_{2q}^\nu}\geq0.}\tag{HC15}
\]
Indeed $m_{h,k}(u)\,du\geq d\nu_k$ bounds every source polynomial
quadratic form.  Minimizing on each identical affine remainder fibre
gives $G_N\succeq G_N^\nu\succ0$.  Conjugation by
$(G_N^\nu)^{-1/2}$ shows that each volume ratio in (HC15) is at least
one.  Cancellation proves the equality.  The low-degree factors require
only original moments through $2q$, and the comparison moments in the
last factors go through $4q$.  These are well-defined finite expressions,
with the actual full jet map and source mass.  The supplied proof
records do not evaluate their errors at the hypothetical arithmetic
quartet or bound them at a rate opposing (HC5).  Retaining this sharper
finite expression therefore supplies no extra contradictory inequality.
